\section{Seed study: results}
\label{app:seed-results}

\begin{table}[t]
\centering
\scriptsize
\setlength{\tabcolsep}{4pt}
% Morandi tints for the cf-transfer tables (muted, low saturation); darker = larger value
\definecolor{cfSage1}{HTML}{EEF1EA}
\definecolor{cfSage2}{HTML}{D9E1D2}
\definecolor{cfSage3}{HTML}{C1CDB8}
\definecolor{cfSage4}{HTML}{A6B79C}
\definecolor{cfSlate1}{HTML}{ECEFF2}
\definecolor{cfSlate2}{HTML}{D3DAE2}
\definecolor{cfSlate3}{HTML}{B7C3D0}
\definecolor{cfSlate4}{HTML}{98A8B9}
\definecolor{cfTerra1}{HTML}{F5ECE9}
\definecolor{cfTerra2}{HTML}{E9D3CC}
\definecolor{cfTerra3}{HTML}{DAB5AA}
\definecolor{cfTerra4}{HTML}{C9968A}
\definecolor{cfOat1}{HTML}{F4EFE6}
\definecolor{cfOat2}{HTML}{E8DDC8}
\definecolor{cfGrey}{HTML}{E4E1DC}

\caption{\textbf{Seed replication: Qwen2.5-VL-7B on NIH ChestX-ray14, 600 new patients.} Per concept: probe selectivity $S$, clean-answer AUROC, concept write $W_{qq}$, the random-family 95th percentile and the absolute sham at the same dose, the rank of $W_{qq}$ among the 120 compared effects, the ownership contrast $O_q$ with its 95\% percentile interval, the strongest competing direction, and the max-$T$ verdict. Negative ownership is shaded terracotta.}
\label{tab:cf-seed}
\begin{tabular}{lrrrrrrrll}
\toprule
Concept & $S$ & AUROC & $W_{qq}$ & p95 & $|$sham$|$ & rank & $O_q$ [95\% CI] & Competitor & Verdict \\
\midrule
Effusion & 0.102 & 0.593 & +0.190 & 0.137 & 0.017 & 2/120 & \cellcolor{cfTerra2}-0.071 [-0.080, -0.062] & Nodule & competitor \\
Atelectasis & 0.062 & 0.566 & +0.101 & 0.297 & 0.096 & 55/120 & \cellcolor{cfTerra3}-0.226 [-0.233, -0.219] & Nodule & competitor \\
Pneumothorax & 0.152 & 0.508 & +0.089 & 0.320 & 0.113 & 52/120 & \cellcolor{cfTerra3}-0.102 [-0.108, -0.096] & Effusion & competitor \\
Cardiomegaly & 0.205 & 0.742 & +0.085 & 0.311 & 0.158 & 33/120 & \cellcolor{cfTerra3}-0.219 [-0.227, -0.211] & Atelectasis & competitor \\
Mass & 0.021 & 0.678 & +0.359 & 0.372 & 0.092 & 10/120 & \cellcolor{cfTerra3}-0.169 [-0.178, -0.160] & Effusion & competitor \\
Nodule & -0.041 & 0.642 & +0.118 & 0.158 & 0.055 & 15/120 & \cellcolor{cfTerra2}-0.072 [-0.081, -0.063] & Effusion & competitor \\
\bottomrule
\end{tabular}
\end{table}


\subsection{A non-generic write effect reproduces}
\label{sec:qwen-replication}

The first finding starts with a successful intervention: Qwen Effusion reproduces beyond random and sham controls before clinical alternatives change its attribution.

Controlled decodability is evaluated at the final visual block each connector consumes: the penultimate CLIP encoder layer for LLaVA and the last of Qwen's 32 visual blocks. Intervals are 95\% patient-cluster bootstrap intervals over 2,000 draws; the control AUROC is the mean over the 20 recurring-type random-label maps, with its 5th--95th percentile spread in parentheses.

Qwen Effusion first satisfies the reader test. Its probe reaches AUROC \cfSeedDecQwenEffAuroc{} $[\cfSeedDecQwenEffAurocLo,\cfSeedDecQwenEffAurocHi]$ against a control AUROC of \cfSeedDecQwenEffCtrl{} (\cfSeedDecQwenEffCtrlLo--\cfSeedDecQwenEffCtrlHi), a controlled selectivity of \cfSeedDecQwenEffSel{} (95\% CI $[\cfSeedDecQwenEffSelLo,\cfSeedDecQwenEffSelHi]$). LLaVA-1.5-7B reads both of its concepts at its consumed block: Effusion at AUROC \cfSeedDecLlavaEffAuroc{} $[\cfSeedDecLlavaEffAurocLo,\cfSeedDecLlavaEffAurocHi]$ with selectivity \cfSeedDecLlavaEffSel{} $[\cfSeedDecLlavaEffSelLo,\cfSeedDecLlavaEffSelHi]$, and Edema at AUROC \cfSeedDecLlavaEdemaAuroc{} $[\cfSeedDecLlavaEdemaAurocLo,\cfSeedDecLlavaEdemaAurocHi]$ with selectivity \cfSeedDecLlavaEdemaSel{} $[\cfSeedDecLlavaEdemaSelLo,\cfSeedDecLlavaEdemaSelHi]$, against a control AUROC of \cfSeedDecLlavaEffCtrl{} (\cfSeedDecLlavaEffCtrlLo--\cfSeedDecLlavaEffCtrlHi). Original-cohort intervention outcomes maximize the concept-consistent effect over six fixed nonzero doses. Qwen makes a strong write-direction case in the original 200-image cohort: the maximum concept-consistent change is 0.2343 at $\alpha=+0.25$, above the random-direction 95th percentile of 0.0963 and the maximum absolute sham effect of 0.0852.

The response at the locked dose $\alpha=+0.25$ reproduces on the patient-disjoint 400-patient cohort. Effusion changes mean $P(\mathrm{yes})$ by 0.1945, again above the random-direction 95th percentile of 0.0878 and absolute sham effect of 0.0157. None of these patients appears in the original intervention cohort, and the dose was not reselected, supporting a large and reproducible non-generic effect within ChestX-ray14.


\subsection{Semantic competition changes attribution, not efficacy}
\label{sec:specificity-results}

The fixed clinical family provides a harder identification test than non-semantic controls. Nodule changes the Effusion answer by 0.2519 and realizes the point-estimate maximum among the five alternatives. The Effusion ownership contrast $O_{\mathrm{Effusion}}$ is $-0.0573$. The registered fixed-family endpoint recomputes the maximum inside each of 5,000 patient-bootstrap draws, giving a two-sided 95\% CI of $[-0.0694,-0.0450]$ and a one-sided 95\% lower bound of $-0.0678$. The entire interval is below zero. Descriptively, all five alternative clinical normals move the Effusion endpoint in the same positive direction; Effusion exceeds four by 0.1172--0.1778, while Nodule alone has a larger point estimate. This shared-sign pattern motivates an answer-sensitive-alias hypothesis but cannot distinguish it from question-specific ownership.

Qwen Effusion therefore remains a reproducible direction that changes the answer beyond random and sham perturbations, while its specific semantic attribution fails. The familywise endpoint establishes that at least one fixed clinical alternative is stronger; Nodule is the point-estimate maximizer, not a separately familywise-identified semantic owner.

\subsection{Wording changes clinical competition within the answer interface}
\label{sec:encoding-results}

The second finding concerns the conditions of clinical competition. The Qwen ownership pattern persists in logit coordinates, so we test whether changing the answer interface changes a direction's relative advantage.

In the registered crossover, Atelectasis, Pneumothorax, Mass, and Nodule pass independent capability calibration. On these four questions, mapping-even minus mapping-odd response energy is $-2.1786$ (95\% CI $[-2.4061,-1.9810]$): effects are larger in the component that reverses raw A-minus-B sign when the finding-to-letter mapping is exchanged. Yet all 20 random directions and the sham are also mapping-odd dominant. Aggregate mapping response therefore does not distinguish the clinical normals from generic perturbations.

The crossover identifies Mass as a prompt-conditioned candidate: its clinical margin changes from $-0.7181$ under the original yes/no prompt to $0.2825$ and $0.2900$ under A-present and B-present, respectively (Table~\ref{tab:seed-mass}, top). A-present assigns A to presence and B to absence; B-present reverses those assignments. Both coded effects exceed the discovery study's 20 random directions and sham. The top rows of Table~\ref{tab:seed-mass} carry pointwise 95\% intervals. Because wording and instructions also change, we test them separately on independent patients.

\begin{table}[H]
\centering
\scriptsize
\setlength{\tabcolsep}{5pt}
\caption{\textbf{Mass clinical competition across prompts.} Clinical margin per prompt: the Mass effect minus the largest of the five other clinical effects, in finding-present logit units. Top, the discovery crossover; bottom, the registered confirmation on new patients.}
\label{tab:seed-mass}
\begin{tabular}{lrrrrr}
\toprule
Prompt & Clinical margin & 95\% CI & Mass effect & Random max & $|$Sham$|$ \\
\midrule
Standard yes/no & $-0.7181$ & $[-0.7956,-0.6412]$ & 2.8625 & 3.0412 & 0.5869 \\
A-present & $0.2825$ & $[0.2269,0.3406]$ & 1.6125 & 1.1650 & 0.7550 \\
B-present & $0.2900$ & $[0.2300,0.3506]$ & 1.7469 & 1.2581 & 1.1475 \\
\midrule
Condition & Clinical margin & Minimum LCB & Mass effect & Random max & \\
\midrule
\textit{show}/A-present & 0.2525 & 0.1675 & 1.7019 & 2.2100 & \\
\textit{show}/B-present & 0.2338 & 0.1462 & 1.8188 & 2.2169 & \\
\textit{is}/A-present & $-0.1087$ & $-0.1751$ & 1.9225 & 2.4750 & \\
\textit{is}/B-present & $-0.2006$ & $-0.2686$ & 2.0019 & 2.4944 & \\
\bottomrule
\end{tabular}
\end{table}


The bottom of Table~\ref{tab:seed-mass} reports the registered confirmation on 200 new patients. Clinical margins subtract the largest of five alternative clinical effects in finding-present logit units. Minimum LCB is the smallest simultaneous lower bound among a cell's five contrasts, from the joint 20-comparison family. All seven prompt cells pass separate capability calibration.

The clinical advantage reproduces under both \textit{show} coded prompts, but the Mass effects remain below the maxima of 119 random directions, with reference ranks of 3/120. These 119 directions form a fixed reference family: each rank places Mass third among the 120 compared effects. This finite descriptive ranking is not a $p$-value, and failure to exceed this family's maximum does not establish nonspecificity across arbitrary direction distributions or reference-family sizes. Under \textit{is}, both clinical margins are negative. Neither the cross-wording nor original-wording specificity conjunction is met.

Explicit \textit{show} yes/no also has a positive clinical margin, 0.2369; its paired differences from A-present and B-present span zero. Within each response instruction, changing \textit{show} to \textit{is} reduces the clinical margin. These paired wording contrasts are descriptive. They support wording-conditioned clinical competition, rather than an advantage attributable to answer letters alone. Mass steering increases Brier error in every cell, while all AUROC-change intervals span zero.

\FloatBarrier
